\begin{frame}<1-3>[t]{新知探索}
\begin{campuscanvas}
% Source slide 7, objects 8; visible 2-3
\only<2-3|handout:1>{\node[anchor=north west,inner sep=0,text=black] at (70.000,190.000) {\begin{minipage}[t]{142.500mm}\raggedright\fontsize{14}{21.00}\selectfont \textbf{例 1}\quad 设 $L(A,B)$ 表示直线 $AB$ 上全体点组成的集合，$P\in L(A,B)$ 的含义是什么？\end{minipage}};}
% Source slide 7, objects 12; visible 3
\only<3|handout:1>{\node[anchor=north west,inner sep=0,text=black] at (70.000,390.000) {\begin{minipage}[t]{142.500mm}\raggedright\fontsize{14}{21.00}\selectfont \red{解}\quad $P\in L(A,B)$ 表示“$P$ 是直线 $AB$ 上的一个点”。\end{minipage}};}
\end{campuscanvas}
\end{frame}
